% TOP_PROBLEM: {"id": 20003307, "title": "Multiplicative minimum and bounded diagonal orbits", "category_id": 20, "category": {"id": 20, "name": "miscellaneous", "display_name": "Miscellaneous", "description": "Problems whose source classification does not fit the main mathematical categories.", "slug": "miscellaneous", "order_index": 20, "created_at": "2026-07-31T15:26:25.671Z"}, "status": "partially_solved"}
% FIRST_POSTED: 2026-09-27
% !TeX program = lualatex
% ATTEMPT_STATUS: unresolved
% Model: GPT-6 Astra Ultra; continuation dated 27 September 2026.
\documentclass[11pt]{article}
\usepackage[margin=25mm]{geometry}
\usepackage{fontspec}
\setmainfont{Latin Modern Roman}
\usepackage{amsmath,amssymb,mathtools}
\usepackage{unicode-math}
\setmathfont{Latin Modern Math}
\usepackage{xurl,hyperref}
\hypersetup{colorlinks=true,urlcolor=blue}
\setcounter{secnumdepth}{0}
\setlength{\parindent}{0pt}
\setlength{\parskip}{5pt}
\setlength{\emergencystretch}{3em}
\begin{document}
\section*{\texorpdfstring{Multiplicative minimum and bounded diagonal orbits}{Multiplicative minimum and bounded diagonal orbits}}\label{20003307}
\subsection{Definitions and mathematical statement}
Let $X=\operatorname{SL}(3,{\mathbb{R}})/\operatorname{SL}(3,{\mathbb{Z}})$, identified with unimodular lattices in ${\mathbb{R}}^3$, and let $A$ be the determinant-one diagonal subgroup. For a lattice $\Lambda$, the orbit $A\Lambda$ is bounded here if its closure in $X$ is compact. Equivalently, there is ${\varepsilon}>0$ such that every nonzero vector of every lattice $a\Lambda$, $a\in A$, has Euclidean length at least ${\varepsilon}$. The conjecture is
\[
 \overline{A\Lambda}\text{ compact}\quad\Longrightarrow\quad
 A\Lambda\text{ closed in }X.
\]
A closed bounded orbit is therefore compact. The statement concerns the whole diagonal-group orbit, not an orbit under one chosen diagonal matrix or one ray of matrices. ``Bounded'' is relative compactness, not boundedness for an arbitrary nonproper metric.
\par\smallskip\textit{Formulation and concept references:} \href{https://arxiv.org/html/2504.17644v3}{[S1]}.

\subsection{Short English statement}
Must every relatively compact orbit of the full diagonal subgroup on the space of three-dimensional unimodular lattices already be closed?
\subsection{Research attempt}
\paragraph{The real rank-two problem is retained.}
The inspected source [S1, Introduction] states the real diagonal-orbit
conjecture separately from its main theorem. That theorem constructs
nonclosed bounded orbits in dimension four over specified positive
characteristic local fields. It is not a counterexample in
$\operatorname{SL}(3,\mathbb R)/\operatorname{SL}(3,\mathbb Z)$.
The present attempt works throughout with real three-dimensional
lattices and the full determinant-one diagonal group.

Let $A^+$ denote the subgroup with positive diagonal entries. Its
logarithm identifies it with
\[
 H=\{(t_1,t_2,t_3)\in\mathbb R^3:t_1+t_2+t_3=0\}.
\]
The full $A$ is a union of four sign translates of $A^+$. Thus
compactness and boundedness statements for $A^+$ immediately pass
to the full group by a finite union. The distinction between this
two-dimensional group and a single diagonal flow remains essential.

\paragraph{A multiplicative reformulation of boundedness.}
For a nonzero vector $v=(v_1,v_2,v_3)$ with all coordinates nonzero,
the arithmetic-geometric mean inequality gives
\[
 \sum_{i=1}^3e^{2t_i}|v_i|^2
       \geq3|v_1v_2v_3|^{2/3}\qquad(t\in H).
\]
Equality is obtained by choosing
$e^{t_i}=|v_1v_2v_3|^{1/3}/|v_i|$, whose product is one.
Consequently
\[
 \inf_{a\in A^+}\|av\|_2=\sqrt3\,|v_1v_2v_3|^{1/3}.
 \tag{393.1}
\]
If a coordinate vanishes, the same infimum is zero: shrink all
nonzero coordinates by $e^{-s}$ and expand the zero coordinates
to restore determinant one. This also agrees with (393.1).

Taking the infimum over nonzero lattice vectors, and using the
compactness criterion already specified in the question, proves
the exact equivalence
\[
 \overline{A\Lambda}\text{ is compact}
 \quad\Longleftrightarrow\quad
 \inf_{0\ne v\in\Lambda}|v_1v_2v_3|>0.
 \tag{393.2}
\]
The exchange of infima is legitimate because both sides are infima
over the same Cartesian product of vectors and group elements.
Signs in $A$ do not change Euclidean lengths or absolute products.
In particular every nonzero vector of a lattice with bounded full
orbit has all three coordinates nonzero. A proposed example containing
a coordinate-axis vector fails the hypothesis immediately.

This criterion also explains why testing a bounded portion of the
diagonal group is insufficient. For any fixed lattice, all diagonal
matrices in a compact parameter set have a positive common least
singular value. Such a finite-parameter test always has a shortest
vector bound, including for lattices whose full orbit escapes.

\paragraph{Periods give a complete sufficient condition for compactness.}
Let
\[
 \Gamma_\Lambda=\{t\in H:
       \operatorname{diag}(e^{t_1},e^{t_2},e^{t_3})\Lambda=\Lambda\}.
\]
This is a discrete additive subgroup. Indeed, if a sequence of such
diagonal transformations tends to the identity, apply it to each
member of a fixed lattice basis. Discreteness of the lattice forces
each basis vector eventually to be fixed, so the transformation is
eventually the identity. The same argument after translation proves
discreteness at any period.

If $\Gamma_\Lambda$ contains two linearly independent elements,
their integer span is a full lattice in the two-dimensional space
$H$. A compact parallelogram represents its quotient. Every $A^+$
translate of $\Lambda$ then has a representative in the continuous
image of that parallelogram. This image is compact and equals
$A^+\Lambda$. The full $A\Lambda$ is a finite union of its sign
translates, hence compact and therefore closed.

Thus a concrete route to the conjecture would be to derive two
independent logarithmic periods from (393.2). The following explicit
arithmetic example shows how this mechanism works when actual units
are available; it does not establish their existence for an arbitrary
lattice with a multiplicative lower bound.

\paragraph{An explicit compact orbit without invoking a unit theorem.}
Consider
\[
 f(X)=X^3-3X+1.
\]
It has three distinct real roots $r_1\in(-2,-1)$,
$r_2\in(0,1)$ and $r_3\in(1,2)$ by sign changes and its degree.
It is irreducible over $\mathbb Q$: the only possible rational roots
are $1,-1$, neither of which is a root. Put $\theta=r_1$ and use
the order $\mathbb Z[\theta]$ with basis $1,\theta,\theta^2$.
Its three real embeddings form the lattice
\[
 \Lambda_0=\{(a+br_1+cr_1^2,
              a+br_2+cr_2^2,
              a+br_3+cr_3^2):a,b,c\in\mathbb Z\}.
\]
The determinant of this embedding matrix is a Vandermonde
determinant. The polynomial discriminant is
$-4(-3)^3-27=81$, so the absolute determinant is nine. Therefore
\[
 \Lambda=9^{-1/3}\Lambda_0
 \tag{393.3}
\]
is unimodular.

For a nonzero $\alpha\in\mathbb Z[\theta]$, multiplication by
$\alpha$ in the integral basis has an integer matrix. Its determinant
is the product of the three embeddings of $\alpha$; it is nonzero
because multiplication by a nonzero field element is invertible.
Hence this product has absolute value at least one. Every nonzero
vector of (393.3) therefore satisfies
\[
 |v_1v_2v_3|\geq1/9.
 \tag{393.4}
\]
By (393.1), all diagonal translates have shortest-vector length at
least $\sqrt3\,9^{-1/3}$. This establishes boundedness with a concrete
constant, rather than from a finite numerical search.

There are two visible units:
\[
 \theta^{-1}=3-\theta^2,\qquad
 (\theta-1)^{-1}=\theta^2+\theta-2.
 \tag{393.5}
\]
Both identities follow by reducing modulo $f(\theta)=0$.
Their squares are totally positive and have field norm one. Thus
the two diagonal matrices
\[
 U=\operatorname{diag}(r_1^2,r_2^2,r_3^2),\qquad
 V=\operatorname{diag}((r_1-1)^2,(r_2-1)^2,(r_3-1)^2)
 \tag{393.6}
\]
have determinant one and preserve $\Lambda$, because multiplication
by either unit permutes the order bijectively.

Their logarithmic vectors are independent. The first coordinate of
both is positive. The third coordinate of $\log U$ is positive,
whereas the third coordinate of $\log V$ is negative. Thus they
cannot be positive scalar multiples; their first coordinates rule
out negative scalar multiples. Neither is zero. They give the two
independent periods required above, proving that (393.3) has a
compact closed full diagonal orbit.

This example requires no assertion that every order has a chosen
unit basis. The required units and their independence are exhibited
and checked directly. The arithmetic mechanism is transparent:
integer nonzero norms give the multiplicative lower bound, while
independent units give return periods in every diagonal direction.

\paragraph{What compact closure does give dynamically.}
Let $Y=\overline{A^+\Lambda}$ be compact. For $R>0$, average the
point mass along a square $Q_R=[-R,R]^2$ in linear coordinates on
$H$:
\[
 \nu_R=\frac1{|Q_R|}\int_{Q_R}\delta_{a(t)\Lambda}\,dt.
\]
Compactness supplies a weakly convergent subsequence of these
probability measures. For any fixed $h\in H$ and continuous $F$
on $Y$, translation of the integration domain gives
\[
 \left|\int F\circ a(h)\,d\nu_R-\int F\,d\nu_R\right|
 \leq\|F\|_\infty\frac{|(Q_R+h)\mathbin\triangle Q_R|}{|Q_R|}
 \longrightarrow0.
\]
Every resulting limit is therefore $A^+$-invariant. This is a
genuine deduction from boundedness, but it does not prove that the
original orbit is closed or that the invariant measure has full
support in its closure.

Even full support and ergodicity would not be a purely topological
substitute for periodicity. An irrational linear flow on a two-torus
has dense nonclosed orbits in a compact space and preserves Haar
probability. This is not a lattice-space counterexample; it shows
that an argument deriving closedness from compactness and invariant
averages alone needs a special rigidity input from the homogeneous
space.

Nor does ambient expansion automatically give positive entropy for
the particular invariant measure. On the compact arithmetic orbit
constructed above, the action is a translation on a quotient of
$H$ by its period lattice. A fixed translation is an isometry in a
flat quotient metric. Its number of mutually separated orbit segments
at a fixed scale is bounded independently of segment length, so its
topological entropy is zero, and hence its measure entropy is zero.
Any measure-classification route that requires positive entropy must
establish that extra hypothesis rather than infer it from diagonal
expansion in ambient directions.

\paragraph{Why one-parameter return data do not fill the gap.}
The logarithmic diagonal group has two independent directions.
Finding a return along one ray produces at most one period. A
quotient of $H$ by one rank-one period still has an unbounded
transverse direction. Numerical recurrences are weaker still: closeness
of $a(t)\Lambda$ to $\Lambda$ does not imply equality of lattices.
The diagonal matrices realizing approximate returns can have unbounded
parameters, so the elementary discreteness argument for periods near
the identity does not turn them into exact periods.

The minimization in (393.1) also treats all lattice vectors, rather
than a preferred basis. A basis may have comfortable nonzero
coordinate products while large integer combinations have products
arbitrarily close to zero. Searching combinations in a bounded box
only provides an upper bound on the infimum in (393.2); it cannot
certify its positivity without a structural argument such as the
integral norm in (393.4).

\paragraph{Diagnostics and outcome.}
The exact diagnostic computes multiplication matrices modulo
$X^3-3X+1$, checks both unit inverses and their norm determinants,
and checks nonzero integer norms for a finite sample of algebraic
integers. Root intervals certify the signs used for independence.
The finite norm sample does not replace the determinant argument
valid for every nonzero algebraic integer.

We have an exact multiplicative reformulation, a sufficient period
criterion, and an explicit nontrivial family member satisfying the
desired conclusion. For a general lattice obeying (393.2), the first
unresolved step is obtaining arithmetic periods or another rigidity
mechanism forcing its orbit closed. The function-field theorem in
[S1] does not supply such a step in the real setting. The selected
conjecture remains unresolved by this attempt.
\subsection{Sources}
\begin{itemize}
\item[S1]Qianlin Huang and Ronggang Shi, Bounded diagonal orbits in homogeneous spaces over function fields, Introduction.
\url{https://arxiv.org/html/2504.17644v3}.
\textit{Formulation source.}
\end{itemize}
\end{document}
